% Options for packages loaded elsewhere
\PassOptionsToPackage{unicode}{hyperref}
\PassOptionsToPackage{hyphens}{url}
\PassOptionsToPackage{dvipsnames,svgnames,x11names}{xcolor}
\documentclass[
  10pt,
  fontset=fandol]{ctexart}
\usepackage{xcolor}
\usepackage[margin=16mm]{geometry}
\usepackage{amsmath,amssymb}
\setcounter{secnumdepth}{-\maxdimen} % remove section numbering
\usepackage{iftex}
\ifPDFTeX
  \usepackage[T1]{fontenc}
  \usepackage[utf8]{inputenc}
  \usepackage{textcomp} % provide euro and other symbols
\else % if luatex or xetex
  \usepackage{unicode-math} % this also loads fontspec
  \defaultfontfeatures{Scale=MatchLowercase}
  \defaultfontfeatures[\rmfamily]{Ligatures=TeX,Scale=1}
\fi
\usepackage{lmodern}
\ifPDFTeX\else
  % xetex/luatex font selection
\fi
% Use upquote if available, for straight quotes in verbatim environments
\IfFileExists{upquote.sty}{\usepackage{upquote}}{}
\IfFileExists{microtype.sty}{% use microtype if available
  \usepackage[]{microtype}
  \UseMicrotypeSet[protrusion]{basicmath} % disable protrusion for tt fonts
}{}
\makeatletter
\@ifundefined{KOMAClassName}{% if non-KOMA class
  \IfFileExists{parskip.sty}{%
    \usepackage{parskip}
  }{% else
    \setlength{\parindent}{0pt}
    \setlength{\parskip}{6pt plus 2pt minus 1pt}}
}{% if KOMA class
  \KOMAoptions{parskip=half}}
\makeatother
\setlength{\emergencystretch}{3em} % prevent overfull lines
\providecommand{\tightlist}{%
  \setlength{\itemsep}{0pt}\setlength{\parskip}{0pt}}
\usepackage{amsmath,amssymb,mathtools,needspace}
\xeCJKDeclareCharClass{CJK}{"2032,"0400->"04FF,"2160->"216B,"2460->"2473,"25A0,"25C7,"25CB,"25CF}
\IfFontExistsTF{Arial Unicode MS}{\xeCJKsetup{AutoFallBack=true}\setCJKfallbackfamilyfont{\CJKrmdefault}{Arial Unicode MS}}{}
\setcounter{secnumdepth}{0}
\setlength{\emergencystretch}{2em}
\allowdisplaybreaks
\usepackage{bookmark}
\IfFileExists{xurl.sty}{\usepackage{xurl}}{} % add URL line breaks if available
\urlstyle{same}
\hypersetup{
  pdftitle={算式的建立},
  colorlinks=true,
  linkcolor={Maroon},
  filecolor={Maroon},
  citecolor={Blue},
  urlcolor={Blue},
  pdfcreator={LaTeX via pandoc}}

\title{算式的建立}
\author{}
\date{}

\begin{document}
\maketitle

\subsubsection{11.4.2　算式的建立}\label{ux7b97ux5f0fux7684ux5efaux7acb}

现在运用消元手续具体建立上述奇偶二分法的算式。

回到一般形式的三对角方程组（11.4.3），利用它的第 \(i-1\) 个方程和第 \(i+1\) 个方程从其第 \(i\) 个方程中消去变元 \(x_{i-1}\) 和 \(x_{i+1}\)（不言而喻，其首尾两个方程需做特殊处理），结果加工得出

\[
\begin{cases}
b_i^{(1)}x_i+c_i^{(1)}x_{i+2}=f_i^{(1)},& i=0,1,\\
a_i^{(1)}x_{i-2}+b_i^{(1)}x_i+c_i^{(1)}x_{i+2}=f_i^{(1)},& i=2,3,\cdots,N-3,\\
a_i^{(1)}x_{i-2}+b_i^{(1)}x_i=f_i^{(1)},& i=N-2,N-1,
\end{cases}
\]

式中

\[
a_i^{(1)}=-a_i a_{i-1}/b_{i-1},\quad i=2,3,\cdots,N-1,
\]

\begin{quote}
校注（agent 补充，CH11-E005）：本页第 1 步的第二个（奇数下标）子系统中，原书依次印有 \(a_3^{(1)}x_0\)、\(c_3^{(1)}x_4\)、\(a_5^{(1)}x_2\)、\(c_5^{(1)}x_6\)、\(a_7^{(1)}x_4\)，已据放大原图照录。按该页明确给出的奇数相关链，其变元依次应为 \(x_1,x_5,x_3,x_7,x_5\)；这里五个偶数下标与奇数子系统不相容，属原书下标疑误。
\end{quote}

\begin{quote}
校注（agent 补充，CH11-E006）：11.4.2 开头将``一般形式的三对角方程组''回引为（11.4.3），已照录。式（11.4.3）是 \(N=8\) 的示例，一般形式在式（11.4.1）；此处疑为原书交叉引用错误。
\end{quote}

\href{../../source-images/pdf-291.jpeg}{原页图像}

\[
b_i^{(1)}=
\begin{cases}
b_i-c_i a_{i+1}/b_{i+1},& i=0,1,\\
b_i-a_i c_{i-1}/b_{i-1}-c_i a_{i+1}/b_{i+1},& i=2,3,\cdots,N-3,\\
b_i-a_i c_{i-1}/b_{i-1},& i=N-2,N-1,
\end{cases}
\]

\[
c_i^{(1)}=-c_i c_{i+1}/b_{i+1},\quad i=0,1,\cdots,N-3,
\]

\[
f_i^{(1)}=
\begin{cases}
f_i-c_i f_{i+1}/b_{i+1},& i=0,1,\\
f_i-a_i f_{i-1}/b_{i-1}-c_i f_{i+1}/b_{i+1},& i=2,3,\cdots,N-3,\\
f_i-a_i f_{i-1}/b_{i-1},& i=N-2,N-1.
\end{cases}
\]

可以看出，上述消元手续的特点在于，它从奇（偶）数编号的方程中消去偶（奇）数编号的变元，从而将下标为奇、偶的变元相互分离开来。也就是说，将所给方程组（11.4.3）加工成下标分别为奇、偶的两个子系统：

\[
\begin{cases}
b_0^{(1)}x_0+c_0^{(1)}x_2=f_0^{(1)},\\
a_{2i}^{(1)}x_{2i-2}+b_{2i}^{(1)}x_{2i}+c_{2i}^{(1)}x_{2i+2}=f_{2i}^{(1)},\quad i=1,2,\cdots,N/2-1,\\
a_{N-2}^{(1)}x_{N-4}+b_{N-2}^{(1)}x_{N-2}=f_{N-2}^{(1)};
\end{cases}
\]

\[
\begin{cases}
b_1^{(1)}x_1+c_1^{(1)}x_3=f_1^{(1)},\\
a_{2i+1}^{(1)}x_{2i-1}+b_{2i+1}^{(1)}x_{2i+1}+c_{2i+1}^{(1)}x_{2i+3}=f_{2i+1}^{(1)},\quad i=1,2,\cdots,N/2-1,\\
a_{N-1}^{(1)}x_{N-3}+b_{N-1}^{(1)}x_{N-1}=f_{N-1}^{(1)}.
\end{cases}
\]

可见，上述消元手续是一项奇偶二分手续。反复施行这种二分手续加工 \(k\) 步，其相关链如式（11.4.2）所示，相应地，所给方程组（11.4.3）被加工成如下形式：

\[
\begin{cases}
b_i^{(k)}x_0+c_i^{(k)}x_{i+2^k}=f_i^{(k)},& i=0,1,\cdots,2^k-1,\\
a_i^{(k)}x_{i-2^k}+b_i^{(k)}x_i+c_i^{(k)}x_{i+2^k}=f_i^{(k)},& i=2^k,2^k+1,\cdots,N-2^k-1,\\
a_i^{(k)}x_{i-2^k}+b_i^{(k)}x_i=f_i^{(k)},& i=N-2^k,N-2^k+1,\cdots,N-1.
\end{cases}
\]

仿照 10.3.2 小节关于一阶线性递推的处理方法，不难导出如下算式：

\[
a_i^{(k)}=-a_i^{(k-1)}a_{i-2^{k-1}}^{(k-1)}/b_{i-2^{k-1}}^{(k-1)},\quad i=2^k,2^k+1,\cdots,N-1,
\]

\[
b_i^{(k)}=
\left\{
\begin{aligned}
&b_i^{(k-1)}-c_i^{(k-1)}a_{i+2^{k-1}}^{(k-1)}/b_{i+2^{k-1}}^{(k-1)},\quad i=0,1,\cdots,2^k-1,\\
&b_i^{(k-1)}-a_i^{(k-1)}c_{i-2^{k-1}}^{(k-1)}/b_{i-2^{k-1}}^{(k-1)}-c_i^{(k-1)}a_{i+2^{k-1}}^{(k-1)}/b_{i+2^{k-1}}^{(k-1)},\\
&\qquad i=2^k,2^k+1,\cdots,N-2^k-1,\\
&b_i^{(k-1)}-a_i^{(k-1)}c_{i-2^{k-1}}^{(k-1)}/b_{i-2^{k-1}}^{(k-1)},\\
&\qquad i=N-2^k,N-2^k+1,\cdots,N-1,
\end{aligned}
\right.
\]

\[
c_i^{(k)}=-c_i^{(k-1)}c_{i+2^{k-1}}^{(k-1)}/b_{i+2^{k-1}}^{(k-1)},\quad i=0,1,\cdots,N-2^k-1,
\]

\begin{quote}
校注（agent 补充，CH11-E006）：本页关于一般 \(N\) 阶方程组的两处回引仍印为（11.4.3），已照录，一般形式实际见式（11.4.1）；``10.3.2 小节关于一阶线性递推''亦按原书照录，该主题与本章 11.3.2 相符。这些属于原书交叉引用疑误。
\end{quote}

\begin{quote}
校注（agent 补充，CH11-E007）：本页 \(b_i^{(1)},f_i^{(1)}\) 的边界分段及后续 \(b_i^{(k)}\)（下一页还有 \(f_i^{(k)}\)）按原书照录，但与所述同时消去左右邻元的手续不一致。例如在分母非零的前提下，消去第 1 行方程中的 \(x_0,x_2\) 应给出 \(b_1^{(1)}=b_1-a_1c_0/b_0-c_1a_2/b_2\) 和 \(f_1^{(1)}=f_1-a_1f_0/b_0-c_1f_2/b_2\)；原书将 \(i=1\) 列入仅含右邻消元项的第一分支，漏去了左邻贡献。\(i=N-2\) 同样漏去右邻贡献。一般 \(k\) 步分段采用新步长 \(2^k\) 作为消元边界，亦需复核；当 \(k=n\) 时首尾分支还发生重叠并出现越界下标。这是原书边界公式疑误，未据此修订正文。
\end{quote}

\begin{quote}
校注（agent 补充，CH11-E008）：两个奇偶子系统的中间行范围均印为 \(i=1,2,\cdots,N/2-1\)，已照录。取上界时分别出现 \(x_N\)、\(x_{N+1}\)，并与单列的末行重复；按该分段写法，中间行的上界疑应为 \(N/2-2\)。
\end{quote}

\begin{quote}
校注（agent 补充，CH11-E009）：二分 \(k\) 步后的分段方程首行印为 \(b_i^{(k)}x_0\)，已据放大原图照录。其范围包含多个 \(i\)，与本节各行的中心变元应为 \(x_i\) 不一致，疑应为 \(b_i^{(k)}x_i\)。
\end{quote}

\href{../../source-images/pdf-292.jpeg}{原页图像}

\[
f_i^{(k)}=
\left\{
\begin{aligned}
&f_i^{(k-1)}-c_i^{(k-1)}f_{i+2^{k-1}}^{(k-1)}/b_{i+2^{k-1}}^{(k-1)},\quad i=0,1,\cdots,2^k-1,\\
&f_i^{(k-1)}-a_i^{(k-1)}f_{i-2^{k-1}}^{(k-1)}/b_{i-2^{k-1}}^{(k-1)}-c_i^{(k-1)}f_{i+2^{k-1}}^{(k-1)}/b_{i+2^{k-1}}^{(k-1)},\\
&\qquad i=2^k,2^k+1,\cdots,N-2^k-1,\\
&f_i^{(k-1)}-a_i^{(k-1)}f_{i-2^{k-1}}^{(k-1)}/b_{i-2^{k-1}}^{(k-1)},\\
&\qquad i=N-2^k,N-2^k+1,\cdots,N-1.
\end{aligned}
\right.
\]

容易看出，按照上述二分消元手续加工 \(n=\log_2 N\) 步，所给方程组最终退化为下列形式：

\[
b_i^{(n)}x_i=f_i^{(n)},\quad i=0,1,\cdots,N-1,
\]

因此有算法 10.5。

\textbf{算法 10.5}　对 \(k=1,2,\cdots\) 直到 \(n=\log_2 N\) 执行上述算式，则

\[
x_i=f_i^{(n)}/b_i^{(n)},\quad i=0,1,\cdots,N-1,
\]

即为方程组（11.4.1）的解。

\begin{center}\rule{0.5\linewidth}{0.5pt}\end{center}

\subsection{相关原页校注（agent 补充）}\label{ux76f8ux5173ux539fux9875ux6821ux6ce8agent-ux8865ux5145}

以下校注来自本节覆盖的原页，可能也涉及同页相邻小节；原文照录与校正说明分开保留。

\subsubsection{原书 PDF 页 292 的校注}\label{ux539fux4e66-pdf-ux9875-292-ux7684ux6821ux6ce8}

\href{../pages/pdf-292.md}{完整原页转录} · \href{../../source-images/pdf-292.jpeg}{原页图像}

校注记录：CH11-E006, CH11-E007。

\begin{quote}
校注（agent 补充，CH11-E006）：本页``因此有算法 10.5''及算法标题均印为 10.5，已照录。按本章算法 11.1---11.4 的连续编号，此算法疑应编号 11.5；属于原书编号疑误。
\end{quote}

\begin{quote}
校注（agent 补充，CH11-E007）：本页 \(f_i^{(k)}\) 的分段定义延续上一页系数更新的边界问题，原式已完整保留。尤其 \(k=n\) 时首、尾范围重叠，公式还可能引用越界的邻元系数；相关直接消元对照见 PDF 291 的校注。
\end{quote}

\subsection{本节覆盖原页的校注记录}\label{ux672cux8282ux8986ux76d6ux539fux9875ux7684ux6821ux6ce8ux8bb0ux5f55}

下列记录按原页关联，含同页相邻小节的校注；计算时同时读取上文校注和完整原页。

\begin{itemize}
\tightlist
\item
  \texttt{CH11-E005}：原书 PDF 页 290
\item
  \texttt{CH11-E006}：原书 PDF 页 290, 291, 292
\item
  \texttt{CH11-E007}：原书 PDF 页 291, 292
\item
  \texttt{CH11-E008}：原书 PDF 页 291
\item
  \texttt{CH11-E009}：原书 PDF 页 291
\end{itemize}

\end{document}
